% ============================================================
% Project : THE PHILOSOPHY OF VIABILITY
% Book    : Book I — Ontology of Viability
% Book ID : VS.B01
% File    : section_07_formal_propositions_and_derived_claims.tex
% Section : Chapter 07, Section 07
% Purpose : To distinguish formal propositions from foundational
%           commitments and axioms, and to establish the conditions
%           under which claims may be regarded as derived,
%           conditional, interpretive, hypothetical, or unresolved.
% ============================================================

\section{Formal Propositions and Derived Claims}
\label{sec:formal-propositions-and-derived-claims}

\subsection{From Foundational Commitments to Propositional Claims}

The formal reconstruction of an ontology requires a distinction
between the commitments that orient the inquiry and the propositions
that articulate more specific consequences of those commitments.
Without this distinction, foundational principles may be treated
as though they were ordinary conclusions, while propositions that
require additional argument may acquire the appearance of
self-evident foundations.

A foundational commitment establishes an adopted orientation
within the ontological inquiry. An axiom serves as a starting
expression within a formal calculus. A proposition, by contrast,
states a claim whose justification or derivation must be examined.
A derived claim is a proposition supported by specified premises
through an identifiable inferential or argumentative procedure.

These roles are related, but they are not interchangeable.
A foundational commitment does not automatically become a formal
axiom. An axiom does not acquire ontological truth merely by being
adopted within a calculus. A proposition does not become a
foundational commitment merely because it follows from a set of
premises.

The central task of this section is to establish a disciplined
account of propositional claims and their justificatory conditions.
The reconstruction must make clear which propositions express
the direct content of adopted commitments, which follow from
explicit formal premises, which depend on additional assumptions,
and which remain interpretive or unresolved.

The governing thesis is that the status of a proposition must
reflect both the form of its claim and the grounds on which it
is supported. A proposition may be formally derivable without
being unconditionally established as an ontological truth.

\subsection{The Functional Role of a Proposition}

A proposition is a declarative claim that can be examined for
its content, scope, consistency, and grounds of support. In a
formal reconstruction, propositions provide the principal
units through which the consequences of definitions, commitments,
and axiomatic formulations can be articulated and evaluated.

A proposition may perform several functions:

\begin{enumerate}
    \item It may articulate a consequence of an adopted
    foundational commitment.

    \item It may express a relation between concepts whose
    meanings have been specified through definitions.

    \item It may state a condition under which another claim
    holds.

    \item It may summarize the result of an argument.

    \item It may formulate a hypothesis for subsequent
    examination.

    \item It may identify an unresolved relation whose
    possible consequences require further analysis.
\end{enumerate}

These functions do not confer the same justificatory status.
The grammatical or symbolic form of a statement cannot determine
whether it is established, provisional, conditional, or open.

For example, the statement

\[
    \varphi
\]

may be presented as a proposition, but its acceptance requires
an account of its premises, definitions, interpretation, and
support. Its appearance as a declarative sentence does not
establish its truth.

Likewise, symbolic notation does not make a claim more
fundamental than an ordinary-language formulation. Formal
representation may improve precision and reveal inferential
structure, but it does not independently supply ontological
justification.

\subsection{Propositions and Foundational Commitments}

The three foundational principles concerning being and identity
provide the primary point of reference for the present
propositional reconstruction:

\begin{enumerate}
    \item Analytical Distinction;
    \item Ontological Inseparability;
    \item Absence of Ontological Priority.
\end{enumerate}

These principles are jointly adopted commitments, but their
content must remain analytically differentiated.

Let:

\[
    AD = \text{Analytical Distinction},
\]

\[
    OI = \text{Ontological Inseparability},
\]

\[
    AP = \text{Absence of Ontological Priority}.
\]

Their joint adoption may be represented as:

\[
    F_{BI} := AD \land OI \land AP.
\]

This notation records the conjunction of the commitments. It
does not establish that one principle can be derived from the
other two, nor does it show that the conjunction is a complete
formal representation of their ontological content.

The distinction between foundational commitment and proposition
can now be stated more precisely:

\begin{itemize}
    \item A foundational commitment establishes an adopted
    ontological orientation.

    \item A proposition articulates a more specific claim
    whose support must be examined.

    \item A derived proposition follows from identified premises
    under specified inferential conditions.

    \item A conditional proposition follows only when additional
    assumptions are granted.
\end{itemize}

The principles therefore guide the formulation and evaluation
of propositions without automatically determining every
proposition that may be advanced within the ontology.

\subsection{The Difference Between Axioms and Propositions}

An axiom is a starting expression accepted within a formal
system. Its role is determined by the structure of that system:
it provides a premise from which further expressions may be
derived according to specified rules.

A proposition, however, is a claim whose justificatory status
may vary. It may be derived from axioms, supported by an
argument, conditionally established, or retained as a hypothesis.

The distinction is functional rather than merely terminological.

Consider a formal system \(\mathcal{L}\) with a set of axioms
\(\mathcal{A}\). A proposition \(\varphi\) may be derivable
from those axioms:

\[
    \mathcal{A} \vdash_{\mathcal{L}} \varphi.
\]

This establishes a formal relation between the axioms and
the proposition. It does not independently establish that
the axioms are ontologically adequate or that the interpretation
of the system corresponds correctly to the intended subject
matter.

An axiom may therefore be formally indispensable while its
ontological justification remains open to examination. A
proposition may be formally derivable while its acceptance
as an ontological claim remains conditional on the adequacy
of the axioms and interpretation.

This distinction prevents the formal system from becoming
a source of unjustified ontological commitments merely through
the internal organization of its expressions.

\subsection{The Scope of Derived Claims}

A derived claim is a proposition whose support depends on
specified premises and an identifiable inferential procedure.

Let \(\Gamma\) represent a set of premises, and let
\(\varphi\) represent a conclusion. A formal derivation may
be expressed as:

\[
    \Gamma \vdash_{\mathcal{L}} \varphi.
\]

The conclusion is derivable within \(\mathcal{L}\) from
\(\Gamma\). Its scope is therefore constrained by:

\begin{enumerate}
    \item the content of the premises;
    \item the meanings assigned to the expressions;
    \item the domain over which the expressions are interpreted;
    \item the inference rules adopted by the system;
    \item any additional assumptions required by the derivation.
\end{enumerate}

A conclusion cannot legitimately exceed these conditions.
If the premises concern only a restricted domain, the conclusion
cannot be generalized to every possible existent without
additional support. If the premises establish only a conditional
relation, the conclusion cannot be presented as unconditional.
If the premises concern a formal representation, the conclusion
cannot automatically be transferred to ontology without an
adequate interpretive argument.

The scope of a derived proposition must therefore be recorded
as part of its content and justification, not treated as an
incidental qualification.

\subsection{Propositional Consequences of Analytical Distinction}

The first foundational principle concerns the analytical
distinction between being and identity.

Its formal expression in the preceding section was represented
through the distinction between the analytical questions or
roles associated with being and identity:

\[
    q_B \neq q_I.
\]

This expression supports a limited but important propositional
consequence.

\paragraph{Proposition 7.1 — Distinct Analytical Questions.}

The analysis of being and the analysis of identity must preserve
their distinct analytical questions or roles.

\paragraph{Justification.}

The proposition articulates the direct content of Analytical
Distinction. If the two analytical questions are treated as
interchangeable, the distinction adopted by the principle is
not preserved.

The proposition is therefore supported by the foundational
commitment, subject to the interpretation of \(q_B\) and
\(q_I\) as analytical questions or roles rather than as
independently existing entities.

\paragraph{Scope.}

The proposition does not establish that being and identity
are separate entities, that they are mutually exclusive, or
that every predicate applicable to one must be inapplicable
to the other.

It also does not establish that their analytical treatment
must remain permanently separate. Their relation may require
joint examination precisely because the distinction operates
within an ontology that also adopts their inseparability.

\paragraph{Proposition 7.2 — Non-Interchangeability of
Analytical Roles.}

A formal or conceptual treatment that substitutes the analytical
role of identity for that of being, or the reverse, requires
an explicit justification.

\paragraph{Justification.}

If the analytical questions are distinct, an unqualified
substitution between their roles would erase the distinction
that the principle requires the analysis to preserve.

The proposition does not prohibit every relation or translation
between the two analytical treatments. It requires that such
relations be justified rather than presumed.

\paragraph{Scope.}

This proposition concerns analytical procedure and conceptual
articulation. It does not establish a complete theory of the
relations between all properties, predicates, or descriptions
associated with being and identity.

\subsection{Propositional Consequences of Ontological
Inseparability}

The second foundational principle concerns the ontological
inseparability of being and identity.

The provisional co-applicability schema introduced in the
preceding section is:

\[
    \forall x \in D_E
    \bigl(E(x) \rightarrow (B(x) \land I(x))\bigr).
\]

Within its stated domain and interpretation, this schema
supports a restricted propositional consequence.

\paragraph{Proposition 7.3 — Co-Applicability within the
Specified Domain.}

For every \(x\) in the specified domain \(D_E\), if \(x\)
satisfies the predicate \(E\), then both \(B(x)\) and \(I(x)\)
hold within the adopted formal representation.

\paragraph{Justification.}

This proposition restates the content of the co-applicability
schema. For an admissible individual \(a\), universal
instantiation yields:

\[
    E(a) \rightarrow (B(a) \land I(a)).
\]

Given \(E(a)\), modus ponens yields:

\[
    B(a) \land I(a).
\]

The inference is valid within the specified formal system.

\paragraph{Scope.}

The proposition establishes co-applicability under the
formalization. It does not independently establish the complete
ontological meaning of inseparability.

In particular, it does not establish:

\[
    B(a) \leftrightarrow I(a),
\]

nor does it establish that being and identity are identical
predicates, interchangeable concepts, or equivalent ontological
determinations.

The broader commitment to ontological inseparability retains
a content that cannot be exhausted by the co-applicability
schema alone.

\paragraph{Proposition 7.4 — Joint Ontological Consideration.}

Within the adopted ontology, the analysis of an existent
must consider being and identity together while preserving
their distinct analytical roles.

\paragraph{Justification.}

This proposition expresses the joint implications of
Analytical Distinction and Ontological Inseparability.

The first principle requires that the analytical questions
remain distinguishable. The second requires that the inquiry
not treat being and identity as ontologically separable.

Their joint adoption therefore supports an approach in which
the two are considered together without collapsing their
analytical roles.

\paragraph{Scope.}

This proposition is an interpretive articulation of the
joint commitments. It is not a formal derivation from the
co-applicability schema alone. Its justification depends
on the adopted ontological interpretation of inseparability
and the intended meaning of analytical distinction.

\subsection{Propositional Consequences of the Absence of
Ontological Priority}

The third foundational principle denies ontological priority
between being and identity in either direction.

Its provisional formal representation is:

\[
    \neg(B \prec_O I)
    \land
    \neg(I \prec_O B).
\]

This expression supports two directly stated consequences.

\paragraph{Proposition 7.5 — Denial of Priority of Being
over Identity.}

Being is not ontologically prior to identity in the sense
specified by the relation \(\prec_O\).

\paragraph{Proposition 7.6 — Denial of Priority of Identity
over Being.}

Identity is not ontologically prior to being in the sense
specified by the relation \(\prec_O\).

\paragraph{Justification.}

Both propositions are obtained through conjunction elimination
from the formal representation of Absence of Ontological
Priority.

Their acceptance depends on the definition and intended
interpretation of \(\prec_O\). The formal expressions cannot
be evaluated independently of that definition.

\paragraph{Scope.}

The propositions deny the two specified directions of
ontological priority. They do not establish the absence of
every possible relation between being and identity.

The following questions remain distinct:

\begin{itemize}
    \item whether one concept is explanatorily prior to another;
    \item whether one description presupposes another;
    \item whether one inquiry is methodologically prior to another;
    \item whether one form of knowledge depends on another;
    \item whether temporal or causal relations obtain.
\end{itemize}

No conclusion concerning these relations follows automatically
from the denial of ontological priority.

\subsection{Joint Propositional Consequences}

The three foundational commitments can be considered together
to formulate a further proposition.

\paragraph{Proposition 7.7 — Distinct but Jointly Required
Ontological Analysis.}

The ontological analysis of being and identity must preserve
their analytical distinction, consider their ontological
inseparability, and avoid assigning ontological priority
to either over the other.

\paragraph{Justification.}

This proposition articulates the combined orientation of
the three foundational commitments:

\[
    AD \land OI \land AP.
\]

Its content follows from their joint adoption as a governing
framework for the analysis. It does not demonstrate that
the principles are mutually derivable or that the proposition
is an exhaustive theory of being and identity.

\paragraph{Scope.}

The proposition governs the present analytical treatment.
It does not settle every possible question concerning the
nature, structure, manifestation, or persistence of identity.

It also does not establish that all ontological distinctions
can be reduced to the relation between being and identity.
Its function is to constrain the reconstruction, not to
close the field of ontological inquiry.

\subsection{Claims That Require Additional Premises}

Several propositions may appear to follow naturally from
the foundational commitments but cannot be derived from
them without additional premises.

The distinction is important because a formally plausible
statement may contain substantive commitments not included
in its stated foundations.

\paragraph{Identity and Determination.}

The foundational principles do not, by themselves, establish
a complete account of the relation between identity and
determination.

The rejection of their automatic equivalence requires
additional conceptual and ontological argument. A proposition
concerning their relation must therefore be supported by
the relevant analysis of determination rather than attributed
to the three principles as a direct formal consequence.

\paragraph{Identity and Cognition.}

The three principles do not, by themselves, provide a complete
formal proof that identity is independent of cognition.

The ontological independence of identity from recognition,
language, awareness, or description is an adopted constraint
of the present ontology. It must be preserved as a distinct
commitment and supported through its own argument.

It cannot be derived merely from the analytical distinction
between being and identity or from the denial of ontological
priority between them.

\paragraph{Identity Plurality.}

The three principles do not determine whether multiple
identities can coexist within one existent, whether one
existent may participate in several identities, or whether
identity plurality requires a particular form of individuation.

These questions remain open. A plurality of components,
relations, manifestations, or descriptions cannot be used
as a substitute for an argument establishing plurality of
identities.

\paragraph{Persistence and Continuity.}

The principles concerning being and identity do not, by
themselves, establish the conditions under which an identity
persists through change or remains continuous across different
manifestations.

Such claims require additional premises concerning the
relevant transformations, respects of comparison, and
conditions of identity.

\paragraph{Unity and Singularity.}

The joint consideration of being and identity does not
establish that every existent is numerically singular or
that unity must be interpreted as indivisibility, simplicity,
or the absence of internal plurality.

These matters require independent conceptual and ontological
examination.

The purpose of identifying these limitations is not to
exclude the corresponding questions. It is to prevent
unjustified conclusions from being attributed to premises
that do not establish them.

\subsection{Four Categories of Propositional Support}

For the present reconstruction, propositions should be
distinguished according to the kind of support available
to them.

\paragraph{Directly Articulated Consequences.}

These propositions express the content of an adopted
foundational commitment with minimal additional interpretation.

Examples include the requirement to preserve the analytical
distinction between being and identity and the denial of
the two specified directions of ontological priority.

Such propositions remain dependent on the meaning and scope
of the commitments they articulate.

\paragraph{Formally Derived Consequences.}

These propositions follow from explicit premises through
specified inference rules.

The co-applicability result in Proposition 7.3 is an example.
Its formal derivation is valid within the adopted schema,
but its ontological significance remains limited by the
interpretation of that schema.

\paragraph{Conditionally Supported Consequences.}

These propositions follow only when additional premises,
definitions, or assumptions are granted.

Their support must remain conditional until those additional
requirements are independently justified.

\paragraph{Interpretive Consequences.}

These propositions articulate the philosophical significance
of a group of commitments or formal results.

Proposition 7.4 and Proposition 7.7 illustrate this category.
They are supported by the intended interpretation of the
foundational principles, but they should not be presented
as purely mechanical deductions from the symbolic schemas.

These four categories concern the manner in which a proposition
is supported. They do not replace the separate epistemic
classification required by the broader claim-status system.

\subsection{The Difference Between Derived and Merely Compatible
Claims}

A claim may be compatible with the foundational commitments
without being derived from them.

Compatibility means that the claim does not appear to
contradict the adopted principles under the relevant
interpretation. Derivability requires more: the conclusion
must follow from specified premises through legitimate
inferential steps.

Suppose a proposition \(\varphi\) is compatible with the
foundational commitments \(F_{BI}\). This does not establish:

\[
    F_{BI} \vdash_{\mathcal{L}} \varphi.
\]

The proposition may require additional premises or an
independent argument.

This distinction prevents a common form of overextension:
a claim is introduced because it appears consistent with
the ontology and is subsequently described as though it
were entailed by the ontology.

Compatibility is a necessary consideration for retaining
a claim within a coherent system, but it is not sufficient
for establishing that claim.

A proposition may therefore be:

\begin{itemize}
    \item compatible but not derived;
    \item derived under additional assumptions;
    \item formally derivable but interpretively inadequate;
    \item inconsistent with an adopted commitment;
    \item unresolved because its relation to the foundations
    has not yet been determined.
\end{itemize}

The reconstruction must preserve these distinctions in both
its exposition and its internal classification of claims.

\subsection{The Status of Propositions Across Revision}

Propositions may change status as the ontology develops.
A claim initially introduced as a hypothesis may later receive
independent support. A formally derived proposition may require
revision when its premises or definitions are refined. A
previously accepted proposition may be rejected if its
inferential structure is shown to be invalid.

Such changes must be recorded without confusing the history
of a claim with its current justificatory status.

For example, if a proposition depends on a definition that
is later revised, the proposition must be re-examined under
the revised definition. It cannot be retained merely because
it was once derived from an earlier formulation.

Likewise, if a proposition is established conditionally,
the later acceptance of its additional premise may strengthen
its status. The conclusion does not become unconditional
retroactively; its present justification must be reconstructed
with the newly adopted premise explicitly included.

This process ensures that formal development remains
revisable without allowing earlier conclusions to persist
after the grounds supporting them have changed.

\subsection{Open Propositions and Unresolved Consequences}

Not every significant ontological claim can be classified
as established or rejected at the present stage.

Some propositions identify genuine possibilities whose
conditions remain undetermined. Others articulate questions
that cannot yet be represented adequately within the available
formal vocabulary.

Such claims should remain open rather than being forced
into an established or rejected category.

The following questions remain unresolved within the
present reconstruction:

\begin{enumerate}
    \item What formal representation can adequately express
    the full ontological content of inseparability without
    collapsing analytical distinction?

    \item What additional premises, if any, are required to
    articulate the relation between identity and determination?

    \item Under what conditions can identity plurality be
    established, and what distinctions are required to
    distinguish it from plurality of components or manifestations?

    \item How should the independence of identity from
    cognition be represented without reducing ontology to
    a theory of predicates or descriptions?

    \item What formal conditions are required for propositions
    concerning persistence and continuity?

    \item Which additional axiomatic candidates may be
    justified without introducing commitments that exceed
    the present ontological foundations?
\end{enumerate}

These questions are not defects in the formal reconstruction.
They identify the boundaries between what has been adopted,
what has been derived, and what still requires investigation.

\subsection{A Protocol for Evaluating Formal Propositions}

Every substantial proposition introduced in the reconstruction
should be examined according to the following requirements:

\begin{enumerate}
    \item \textbf{Statement:} formulate the proposition clearly
    and identify its intended subject matter.

    \item \textbf{Classification:} determine whether it is a
    foundational articulation, a formally derived proposition,
    a conditional consequence, an interpretive claim, or a
    hypothesis.

    \item \textbf{Premises:} identify the commitments,
    definitions, axioms, and additional assumptions on which
    it depends.

    \item \textbf{Derivation:} where a formal derivation is
    claimed, display the relevant inferential steps.

    \item \textbf{Interpretation:} specify the intended meanings
    and domain of the expressions used.

    \item \textbf{Scope:} state the conditions and limits of
    the proposition.

    \item \textbf{Status:} assign an epistemic or justificatory
    status independently of its functional form.

    \item \textbf{Consistency:} examine its relation to the
    foundational commitments and previously established claims.

    \item \textbf{Revision conditions:} identify what changes
    in premises, definitions, or interpretation would require
    the proposition to be reconsidered.
\end{enumerate}

This protocol provides a basis for the later construction
of a complete propositional and dependency register. It also
ensures that formal precision does not conceal unresolved
philosophical assumptions.

\subsection{Concluding Principle}

The role of propositions in the formal reconstruction is
to articulate the consequences, conditions, and limits of
the adopted ontological commitments.

A proposition must not be treated as established merely
because it is compatible with the foundations, expressed
in symbolic notation, or presented as a consequence of
a respected principle. Its support must be traceable to
identified premises and legitimate inferential or
interpretive procedures.

The governing principle may therefore be stated as follows:

\begin{quote}
A formal proposition acquires the status of a derived claim
only to the extent that its premises, inferential conditions,
interpretation, and scope are explicitly established. Claims
that require additional assumptions must remain conditional,
while claims whose support is unresolved must remain open.
\end{quote}

This principle preserves the distinction between foundational
commitments, axioms, propositions, and derived claims. It
also provides the necessary basis for examining the arguments
and proofs through which the propositions of the reconstruction
are to be supported.

% ============================================================
% END OF SECTION 07 OF CHAPTER 07
% ============================================================